% !TEX root = ../main.tex
\section*{Executive Summary}
\addcontentsline{toc}{section}{Executive Summary}
\label{sec:executive-summary}

\noindent\textbf{致：刚果民主共和国卫生部及国际捐助方}

\medskip
按题面统计，截至 2026 年 10 月 25--26 日，本轮疫情累计确诊约 10\,060 例、死亡约 4\,887 人，已波及 8 个省，肯尼亚与尼日利亚出现输入病例。社区发现不足、医疗服务受冲突影响和资金短缺共同制约疫情应对。

\paragraph{主要发现。}
研究资料延续至 11 月 17 日，期间新增确诊合计为 11\,605 例。参考病例预测能描述历史走势，但后七周回测 RMSE 为 119.50 例，高于持续性基准的 60.10 例；预测须随新增资料滚动更新。省际人口守恒并不保证传播总量不变，地区间的发现和救治能力应与流动情况同时评估。

\paragraph{近期行动。}
将病例发现、规范收治、院内防护、安全安葬与接触者随访作为连续服务组织实施。优先保障一线防护物资、转诊交通和安葬服务的可达性；对受冲突影响的地区，与当地社区协调服务安排和人员安全。追踪质量以实际随访、症状发现至转诊的时间及失访情况评价。

\paragraph{资源与资金。}
若以 2.1 万人为等效随访人力基数且人均工作量不变，随访率从 78\% 提至 90\% 或 95\%，对应新增约 3\,231 人或 4\,577 人。预算须结合专职投入、工作周期、交通与安全成本核定。题面所列超过 4 亿美元的应对资金缺口，应按可核验的人员、物资和服务项目分配，并定期公布资金执行与服务覆盖情况。

\paragraph{区域协作。}
建立跨境病例信息与转诊协作机制，及时共享病例关联和接触者调查结果。肯尼亚和尼日利亚的响应依据当地传播条件调整，持续监测输入后的病例发现和服务能力。针对本迪布焦病毒的疫苗及特异性治疗，按临床试验和监管进展更新应用安排。

\paragraph{对国家安全的投资。}
当地救治和监测能力的提升，有助于减少跨境输入、医疗挤兑及后续应急支出，并维持人员往来和基本公共服务。国际支持的收益超出单一国家，因此稳定的区域防控投入既具有公共卫生价值，也关系到周边国家的社会与经济安全。
